\section{Methods}

A focal active-inference agent plays a repeated trust game with four partners whose behavioral types and stances are hidden and must be inferred from behavioral evidence. The agent is implemented using the \texttt{pymdp} library for discrete active inference \cite{Heins2022Pymdp,DaCosta2020DiscreteSynthesis}, which provides standard components for belief updating and policy selection under the active inference framework \cite{Friston2010FreeEnergy,Friston2017ProcessTheory,Parr2022ActiveInference}. Partner-local affective precision is layered on top of per-partner generative models, allowing evidence and action confidence to remain relationship-specific. Full POMDP specifications, affective precision update rules, and simulation protocols are provided in Appendix~\ref{app:pomdp}, Appendix~\ref{app:precision_update}, and Appendix~\ref{app:protocols}, respectively.

\subsection{Trust game}

The trust game \cite{Berg1995Trust,KingCasas2005Trust} provides a natural testbed because partner types are hidden and must be inferred, creating a genuine gap between social knowledge and actionable confidence. Each partner has a latent \emph{type}---cooperator, reciprocator, exploiter, or random---and a latent \emph{stance} toward the agent---trusting, neutral, or hostile---neither directly observable; both must be inferred from behavioral evidence. Episodes run for up to 200 rounds.

In the binary game both players choose to cooperate or defect, with focal agent payoffs as shown in Table~\ref{tab:payoff}. The structure poses a social inference problem: defection is locally tempting, but cooperative partners are worth identifying and maintaining over repeated interaction. In the graded regime, the binary choice is replaced by a discretized investment level. In the partner-choice regime, the agent selects which partner to engage before choosing an action each round, making approach and avoidance explicit.

\begin{table}[H]
\centering
\caption{Focal agent payoffs in the binary trust game.}
\label{tab:payoff}
\begin{tabular}{lcc}
\toprule
 & Partner cooperates & Partner defects \\
\midrule
Agent cooperates & $3$ & $-1$ \\
Agent defects & $5$ & $1$ \\
\bottomrule
\end{tabular}
\end{table}

Reported experiments use one focal active-inference agent interacting with four environment-side partners. Partners sample actions from type-by-stance cooperation probabilities and undergo action-dependent stance transitions, but do not perform variational inference or maintain affective precision estimates. This isolates the focal agent's precision mechanism before extending to reciprocal active-inference partners.

\subsection{Per-partner generative models}

Rather than maintaining a single pooled social model, the focal agent maintains a separate discrete generative model $\mathcal{M}_k$ for each partner $k$. This factorization keeps evidence about each partner independent and allows action confidence to vary across relationships. Each model encodes observation likelihoods, transition dynamics, preferences, priors, and policy priors using the standard $A/B/C/D/E$ structure of discrete active inference \cite{DaCosta2020DiscreteSynthesis}. Full matrix specifications are provided in Appendix~\ref{app:pomdp}.

Each partner model contains three latent factors: partner type, partner stance, and a deterministic own-action bookkeeping factor that makes the agent's executed action available to the payoff likelihood. Observations are partner action and focal payoff. The partner-action likelihood is defined by type- and stance-conditioned cooperation probabilities, operationalizing predictive reliability: both a reliably cooperative and a reliably defecting partner can be predictable, even though their payoff consequences differ. Payoff enters as a second observation modality, with partner action marginalized through the cooperation likelihood, so preferences over outcomes remain part of the generative model rather than an external reward cache. Stance transitions are action-dependent---cooperation tends to build trust, defection damages it, and recovery from hostility is slower than collapse.

\subsection{Partner-local affective precision}

Hesp et al.\ \cite{Hesp2021DeeplyFeltAffect} treat policy precision as an inferred quantity, modeling affect as active inference over a higher-order belief $\beta$ representing subjective model fitness. We extend this to the per-partner case by assigning each partner its own $\beta_k$ estimate, updated from that partner's behavioral evidence alone.

After interacting with partner $k$, the agent computes surprisal from the predicted probability of the observed partner response:
\begin{equation}
\epsilon_k = -\log P\!\left(o_k^{\text{act}} = \hat{o}_k^{\text{act}} \;\middle|\; q(s_k^{\text{type}}, s_k^{\text{stance}})\right).
\label{eq:epsilon}
\end{equation}
This is computed from partner action rather than payoff because action is the most direct test of type-by-stance model fit: a reliably defecting partner produces no action surprise even though the payoff is poor. The neutral fitness baseline is set at $\sigma_0^2 = (-\log 0.5)^2$, so a fifty-fifty binary prediction carries zero affective charge. Affective charge is then:
\begin{equation}
\phi_k = \alpha\!\left(\sigma_0^2 - \epsilon_k^2\right),
\label{eq:charge}
\end{equation}
where $\alpha$ is a gain parameter. Positive $\phi_k$ means partner $k$'s model predicted better than baseline; negative $\phi_k$ means it predicted worse. A categorical posterior $q(\beta_k)$ is updated each round by combining a slowly evolving persistence prior with a charge-based pseudo-likelihood; full update rules are given in Appendix~\ref{app:precision_update}. The agent then sets policy precision from the posterior mean:
\begin{equation}
\gamma_k = \gamma_{\text{base}} \,/\, \mathbb{E}_{q(\beta_k)}[\beta_k],
\label{eq:gamma_update}
\end{equation}
where $\gamma_{\text{base}}$ is the baseline policy precision applied in the absence of affective modulation. Here $\beta_k$ represents inverse precision: higher $\beta_k$ means the model has been less reliable, producing a broader, less committed policy distribution. The mechanism thus translates accumulated evidence about predictive reliability into action confidence independently for each social relationship. The $\beta_k$ posterior is maintained as an auxiliary precision tracker rather than a proper hidden state in the generative model, which simplifies implementation at the cost of losing the agent's ability to form prior expectations over its own future affective states; the implications of this design choice are discussed in Section~\ref{sec:discussion}.

\subsection{Cross-partner policy selection}

In partner-choice regimes, policies specify both which partner to engage and how to act toward that partner. A naive application of partner-local precision creates an artifact: low precision toward a uniformly poor partner branch scales its logits toward zero, making it spuriously attractive in cross-partner comparison.

We address this by applying precision to centered branch logits:
\begin{equation}
\tilde{u}_{k,\pi} = \bar{u}_k + \gamma_k\left(u_{k,\pi} - \bar{u}_k\right),
\label{eq:centered_logits}
\end{equation}
where $u_{k,\pi}$ is the pre-precision policy logit for policy $\pi$ toward partner $k$ and $\bar{u}_k$ is that branch's mean logit. Centering preserves the mean score of each partner branch while allowing $\gamma_k$ to sharpen or flatten the within-partner policy distribution. The agent selects from the combined set of centered logits across all partner branches. Full details are provided in Appendix~\ref{app:precision_update}.

\subsection{Simulation setup}

Simulations compare four affect-runtime variants: full partner-local precision, a shared-$\beta$ ablation that pools evidence across partners while keeping per-partner POMDP beliefs intact, a tracked-only ablation that updates $q(\beta_k)$ but fixes $\gamma_k = \gamma_{\text{base}}$, and a no-affect control that disables the $\beta$ tracker entirely. Conditions span binary and graded decision regimes, partner-choice settings, abrupt betrayal, and profile-style parameter variation across precision gain $\alpha$ and prior model fitness. Central comparisons use 30 seeds per condition; profile-scale experiments use 20 (Appendix~\ref{app:protocols}). Full condition definitions, episode lengths, and metrics are provided in Appendix~\ref{app:protocols}. Simulation code, analysis scripts, and experiment configurations will be released in a public repository upon acceptance.
